% !PRINTED 191
\weberpage{191}
%<<CONT

\weq{14}{(\alpha'_h\beta'_k) = r(\alpha_h\beta_k),}

wonach die Determinanten $(\alpha_h\beta_k)$ als irrationale \gesperrt{Invarianten} zu betrachten sind. Wir wollen nun darlegen, wie man daraus rationale Invarianten und Covarianten bilden kann.

Wir betrachten zu diesem Zweck am besten die lineare Transformation in der nicht homogenen Gestalt

\weq{15}{x = \frac{\alpha x' + \beta}{\gamma x' + \delta}}

und wenden diese auf die Function $f(x)$ in (1) an. Es folgt dann, wenn

\[
F(x') = a'_0 x'^n + a'_1 x'^{n-1} + a'_2 x'^{n-2} + \cdots + a'_n
\]

gesetzt wird

\weq{16}{F(x') = (\gamma x' + \delta)^n f(x).}

Diese Gleichung schreiben wir auch so

% [sic] Im Druck fehlt zwischen dem Glied a_2/x^2 und den Auslassungspunkten das sonst erforderliche Pluszeichen.
\[
F(x') = (\alpha x' + \beta)^n \left(a_0 + a_1 \frac{1}{x} + a_2 \frac{1}{x^2} \cdots\right),
\]

woraus, wenn man $x' = -\delta : \gamma$, also $x = \infty$ setzt,

\[
a_0 r^n = (-\gamma)^n F\left(-\frac{\delta}{\gamma}\right),
\]

oder, wenn $\alpha'_1, \alpha'_2, \ldots \alpha'_n$ ebenso von $\alpha_1, \alpha_2, \ldots \alpha_n$ abhängen, wie $x'$ von $x$, also

\[
F(x') = a'_0 (x' - \alpha'_1)(x' - \alpha'_2) \ldots (x' - \alpha'_n)
\]

gesetzt wird,

\weq{17}{r^n a_0 = a'_0(\gamma\alpha'_1 + \delta)(\gamma\alpha'_2 + \delta) \ldots (\gamma\alpha'_n + \delta)}

folgt.

Nun ist nach der Transformation (15)

\weq{18}{\begin{aligned}
x - \alpha_i &= \frac{r(x' - \alpha'_i)}{(\gamma x' + \delta)(\gamma\alpha'_i + \delta)} \\
\alpha_i - \alpha_k &= \frac{r(\alpha'_i - \alpha'_k)}{(\gamma\alpha'_i + \delta)(\gamma\alpha'_k + \delta)}.
\end{aligned}}

Wir bilden nun ein Product $P$ aus irgend welchen der Factoren

\[
\begin{aligned}
&(0, 1), (0, 2) \ldots (0, n) \\
&(1, 2), (1, 3) \ldots (n - 1, n),
\end{aligned}
\]

jedoch so, dass darin jede der Ziffern $1, 2 \ldots n$ im Ganzen gleich oft, etwa $\mu$ mal vorkommt. Der Index $0$, d. h. die Variable $x$, möge $\nu$ mal vorkommen, und die Gesammtanzahl der
%CONT>>
